\clearpage
\chapter*{Annexure: Graduate Map Program Outcome (PO) with Skill Achievement from Industrial Training}
\addcontentsline{toc}{chapter}{Annexure: PO-Skill Mapping}

\par Program Outcomes (POs) \cite{16} serve as critical benchmarks to measure the knowledge, skills, and competencies acquired during my B.Sc. in CSE journey, ensuring alignment between academic curricula and industry demands. While universities focus on theoretical knowledge, Program Outcomes (POs) ensure graduates possess the technical, problem-solving, and ethical skills needed in the industry. During my industrial training, I had the opportunity to apply Computer Science \& Engineering knowledge in real-world scenarios, such as testing, debugging, collaboration, and adhering to ethical standards. This experience bridged the gap between theory and practice, enhanced my technical proficiency, and strengthened my soft skills.

\par Mapping POs is essential to evaluate how well academic learning translates to industry practice. Three core reasons for mapping these outcomes include:
\begin{itemize}
    \item To assess curriculum relevance to industrial expectations,
    \item To evaluate readiness in applying theory to practical tools and technologies, and
    \item To enable self-assessment of professional and ethical growth.
\end{itemize}

\par The following Table~\ref{tab:po-fulfillment} maps the 12 Program Outcomes (POs) with the knowledge and skills acquired through both academic education and industrial training. Fulfillment level indicates how effectively each PO was addressed during the internship.

\begin{longtable}[H]{|p{2.5cm}|p{3.5cm}|p{2cm}|p{4cm}|p{4cm}|}
\caption{Program Outcome (PO) Fulfillment in Industrial Training} \label{tab:po-fulfillment} \\
\hline
\textbf{PO} & \textbf{Description} & \textbf{Fulfillment Level} & \textbf{Training Knowledge Applied} & \textbf{Evidence} \\
\hline
\endfirsthead

\multicolumn{5}{c}%
{{\bfseries \tablename\ \thetable{} -- continued from previous page}} \\
\hline
\textbf{PO} & \textbf{Description} & \textbf{Fulfillment Level} & \textbf{Training Knowledge Applied} & \textbf{Evidence} \\
\hline
\endhead

\hline \multicolumn{5}{r}{{Continued on next page}} \\
\endfoot

\hline
\endlastfoot

\textbf{PO2: Problem Analysis} & Identify, formulate, and analyze complex computing problems. & Fully Fulfilled & Diagnosed root causes of bugs and performance issues in software systems. & Identified a bottleneck in backend queries that improved application speed by 30\%. \\
\hline

\textbf{PO5: Modern Tools} & Apply theoretical understanding of modern tools. & \textit{[Insert Fulfillment Level]} & \textit{[Mention tools used: Git, Postman, JMeter, etc.]} & \textit{[Describe a task or tool automation achieved]} \\
\hline

\textbf{PO6: Engineer \& Society} & Understand societal implications of computing solutions. & \textit{[Insert Fulfillment Level]} & \textit{[Mention any user-data, privacy, or accessibility work]} & \textit{[Evidence of ensuring societal responsibility]} \\
\hline

\textbf{PO7: Environment} & Evaluate environmental impact of computational systems. & Not Fulfilled & N/A & Environmental aspects were not addressed in the assigned project. \\
\hline

\textbf{PO8: Ethics} & Apply ethical theories in decision-making. & Fully Fulfilled & Ensured accuracy in documentation, refused unethical shortcuts. & Reported a misconfiguration instead of bypassing it for faster testing. \\
\hline

\textbf{PO10: Communication} & Communicate technical concepts clearly. & \textit{[Insert Fulfillment Level]} & \textit{[Describe presentation, reporting, or documentation]} & \textit{[Mention audience and feedback or result]} \\
\hline

\textbf{PO12: Lifelong Learning} & Recognize the need for continuous learning. & Fully Fulfilled & Explored a new technology (e.g., Docker, Firebase, or AI testing). & Completed an online course and implemented learned skills in a demo module. \\
\end{longtable}

\begin{longtable}[H]{|p{1.5cm}|p{4cm}|p{4cm}|p{4cm}|}
\caption{Complex Engineering Problems (WP) -- Application to Project}
\label{tab:cep} \\
\hline
\textbf{No.} & \textbf{Attribute} & \textbf{Description} & \textbf{How it applies to your problem} \\
\hline
\endfirsthead
\multicolumn{4}{c}{{\bfseries \tablename\ \thetable{} -- continued from previous page}} \\
\hline
\textbf{No.} & \textbf{Attribute} & \textbf{Description} & \textbf{How it applies to your problem} \\
\hline
\endhead
\hline \multicolumn{4}{r}{{Continued on next page}} \\
\endfoot
\hline
\endlastfoot

\multicolumn{4}{|p{13.5cm}|}{%
  \textbf{Complex problems} have characteristic WP1 and some or all of WP2 to WP7.} \\
\hline

%% --- WP1 (mandatory) ---
\textbf{WP1} &
Depth of Knowledge Required &
Cannot be resolved without in-depth engineering knowledge at WK3, WK4, WK5, WK6 or WK8 level, requiring a first-principles analytical approach. &
\toFill{describe what deep knowledge was needed to solve your problem} \\
\hline

%% --- WP2 ---
\textbf{WP2} &
Range of conflicting requirements &
Involve wide-ranging or conflicting technical, engineering and other issues. &
\toFill{describe any conflicting technical or stakeholder requirements faced} \\
\hline

%% --- WP3 ---
\textbf{WP3} &
Depth of analysis required &
Have no obvious solution and require abstract thinking and originality to formulate suitable models. &
\toFill{explain what analysis or modelling approach was needed} \\
\hline

%% --- WP4 ---
\textbf{WP4} &
Familiarity of issues &
Involve infrequently encountered issues. &
\toFill{describe any unusual or unfamiliar issues encountered} \\
\hline

%% --- WP5 ---
\textbf{WP5} &
Extent of applicable codes &
Are outside problems encompassed by standards and codes of practice for professional engineering. &
\toFill{note if the problem went beyond standard codes or guidelines} \\
\hline

%% --- WP6 ---
\textbf{WP6} &
Extent of stakeholder involvement and level of conflicting requirements &
Involve diverse groups of stakeholders with widely varying needs. &
\toFill{list the stakeholders involved and their differing needs} \\
\hline

%% --- WP7 ---
\textbf{WP7} &
Interdependence &
Are high level problems including many component parts or sub-problems. &
\toFill{describe the sub-problems or components that made up the overall problem} \\
\hline

\end{longtable}

\newpage

\begin{longtable}[H]{|p{1.5cm}|p{4cm}|p{4cm}|p{4cm}|}
\caption{Complex Engineering Activities (EA) -- Application to Project}
\label{tab:cea} \\
\hline
\textbf{No.} & \textbf{Attribute} & \textbf{Description} & \textbf{How it applies to your activity} \\
\hline
\endfirsthead
\multicolumn{4}{c}{{\bfseries \tablename\ \thetable{} -- continued from previous page}} \\
\hline
\textbf{No.} & \textbf{Attribute} & \textbf{Description} & \textbf{How it applies to your activity} \\
\hline
\endhead
\hline \multicolumn{4}{r}{{Continued on next page}} \\
\endfoot
\hline
\endlastfoot

\multicolumn{4}{|p{13.5cm}|}{%
  \textbf{Complex activities} mean (engineering) activities or projects that have some or all of the following characteristics.} \\
\hline

%% --- EA1 ---
\textbf{EA1} &
Range of resources &
Involve the use of diverse resources including people, money, equipment, materials, information and technologies. &
\toFill{list the resources used, e.g.\ team members, tools, budget, datasets} \\
\hline

%% --- EA2 ---
\textbf{EA2} &
Level of interactions &
Require resolution of significant problems arising from interactions between wide-ranging or conflicting technical, engineering or other issues. &
\toFill{describe interactions or conflicts between technical components or teams} \\
\hline

%% --- EA3 ---
\textbf{EA3} &
Innovation &
Involve creative use of engineering principles and research-based knowledge in novel ways. &
\toFill{explain any novel or creative approach taken during the activity} \\
\hline

%% --- EA4 ---
\textbf{EA4} &
Consequences to society and the environment &
Have significant consequences in a range of contexts, characterised by difficulty of prediction and mitigation. &
\toFill{describe potential societal or environmental impact of the activity} \\
\hline

%% --- EA5 ---
\textbf{EA5} &
Familiarity &
Can extend beyond previous experiences by applying principles-based approaches. &
\toFill{describe what was new or unfamiliar and how principles were applied} \\
\hline

\end{longtable}